\documentclass[11pt,a4paper]{article}
\usepackage[T1]{fontenc}
\usepackage[utf8]{inputenc}
\usepackage{lmodern}
\usepackage[english]{babel}
\usepackage{geometry}
\geometry{margin=2.4cm}
\usepackage{graphicx}
\usepackage{booktabs}
\usepackage{longtable}
\usepackage{hyperref}
\usepackage[table]{xcolor}
\usepackage{float}
\usepackage{amsmath}

\newcommand{\repoRoot}{..}
\setlength{\emergencystretch}{3em}
\IfFileExists{generated/results_macros.tex}{\newcommand{\resnRuns}{45}
\newcommand{\resbestAllLabel}{GPT-2 LoRA}
\newcommand{\resbestAllAcc}{70.67}
\newcommand{\resbestAllFone}{70.49}
\newcommand{\resbestPtLabel}{pt-110M full}
\newcommand{\resbestPtAcc}{63.95}
\newcommand{\resbestPtFone}{63.03}
\newcommand{\resbestGptLabel}{GPT-2 LoRA}
\newcommand{\resbestGptAcc}{70.67}
\newcommand{\resgapPtGpt}{6.7}
\newcommand{\resscratchAcc}{57.73}
\newcommand{\resptFiveFullAcc}{62.66}
\newcommand{\resgainPretrain}{4.9}
}{}

\newcommand{\maybeinput}[1]{%
  \IfFileExists{#1}{\input{#1}}{%
    \begin{center}
    \fcolorbox{orange}{yellow!10}{%
      \parbox{0.9\linewidth}{\centering\textit{Plik \texttt{\detokenize{#1}} --- uzupełni się po wygenerowaniu (nocna kolejka).}}}%
    \end{center}}}

\newcommand{\maybefig}[2][]{%
  \IfFileExists{#2}{\includegraphics[width=\linewidth,height=0.82\textheight,keepaspectratio]{#2}}{%
    \begin{center}
    \fcolorbox{orange}{yellow!10}{%
      \parbox{0.85\linewidth}{\centering\textit{Wykres \texttt{\detokenize{#2}} --- pojawi się po zakończeniu eksperymentów. #1}}}%
    \end{center}}}

\title{Fine-tuning klasyfikacji sentymentu (Tweet Eval)\\%
  Porównanie pretrenowanych miniLLaMA i GPT-2}
\author{UMISI --- Computational Linguistics (lab03)}
\date{\today}

\begin{document}
\maketitle
\tableofcontents
\newpage

\section{Cel}
Celem jest porównanie metod dostrajania modeli językowych decoder-only do klasyfikacji
trójklasowej (negatywny / neutralny / pozytywny) na zbiorze \texttt{cardiffnlp/tweet\_eval}
(sentiment). Bazą są własne modele \texttt{transformer\_miniLLaMA} (5--110\,M parametrów)
wytrenowane na WikiText-103 oraz model referencyjny \texttt{openai-community/gpt2}.
Dla każdego backbone'u badamy: pełny fine-tuning (\texttt{full}), LoRA, trening samej głowy
(\texttt{head\_only}) oraz odmrożenie ostatnich $k$ bloków transformera ($k=\lceil pL\rceil$ dla
$p\in\{10\%,20\%,30\%,50\%\}$). Trening klasyfikacji: 5 epok, global batch 64 (2$\times$GPU).

\section{Środowisko obliczeniowe}
\begin{itemize}
  \item Sprzęt: 2$\times$ NVIDIA GeForce RTX 3060 (12\,GiB), AMD Ryzen 5 5600X, 16\,GiB RAM, WSL2 z limitem pamięci 9\,GiB.
  \item Oprogramowanie: Docker, \texttt{uv}, PyTorch + Lightning (\texttt{ddp\_spawn}), \texttt{transformers}, \texttt{peft}, \texttt{datasets}.
  \item Precyzja treningu klasyfikacji: \texttt{bf16-mixed}.
  \item Kod: \texttt{lab03/src/finetune.py}, presety \texttt{finetune\_presets.py}, kolejka \texttt{FINETUNE\_QUEUE}.
\end{itemize}

\section{Pretrening własnych modeli językowych}
\subsection{Korpus i tokenizacja}
\maybeinput{generated/pretrain_prose.tex}

\subsection{Architektura i rozmiary}
\maybeinput{generated/pretrain_table.tex}

\subsection{Uwagi do porównań}
Checkpointy użyte w fine-tunie są przypięte jawnie w \texttt{PRETRAINED\_LM} (pełny WikiText-103).
Wcześniejsze eksperymenty z \texttt{pt-latest} mogły wybierać starsze, krótsze runy smoke --- ta seria tego unika.

\section{Zbiór danych klasyfikacji (Tweet Eval)}
Klasyfikacja pełnotekstowa (nie token-level). Etykiety: \texttt{negative}, \texttt{neutral}, \texttt{positive}.
Dla modeli BPE stosowany jest tokenizer pretreningu; dla GPT-2 --- tokenizer Hugging Face.

\maybeinput{generated/dataset_table.tex}

\maybeinput{generated/dataset_examples.tex}

\textbf{Przetwarzanie wstępne.} Teksty bez zmian treści (anonimizacja \texttt{@user} pochodzi z samego zbioru);
tokenizacja tokenizerem danego backbone'u, ucięcie do 128 tokenów (żaden tweet z próbki go nie przekracza),
dopełnienie do najdłuższej sekwencji w batchu z maską uwagi. Splity są oficjalne, bez re-próbkowania.

\section{Modele klasyfikacyjne}
\textbf{miniLLaMA + głowa:} backbone z checkpointu LM (lub scratch), pooling ostatniego tokena,
MLP (\texttt{head\_hidden}=256, dropout 0.1). \textbf{GPT-2:} \texttt{GPT2ForSequenceClassification}.
Metryki walidacyjne: makro F1; wybór \texttt{best\_eval} po \texttt{val/f1\_macro}.

\section{Siatka eksperymentów (fine-tuning)}
\maybeinput{generated/grid_prose.tex}

\maybeinput{generated/grid_table.tex}

\section{Protokół ewaluacji}
\begin{itemize}
  \item Bez early stoppingu: każdy run trwa pełne 5 epok, po każdej zapisywany jest checkpoint wag
        (\texttt{checkpoints/epoch\_XXXX.pt}). Do testu wybierane są wagi z epoki o najwyższym val F1 macro
        (\texttt{*\_best\_eval.pt}).
  \item Metryki: accuracy, F1 macro i ważone, macierz pomyłek, czas treningu (łącznie i na epokę).
  \item Inferencja: \emph{przepustowość} --- cały test (12\,284 przykłady) batchami po 32 na jednej karcie;
        \emph{opóźnienie} --- 1\,000 pojedynczych przykładów (batch 1), średni czas na przykład.
  \item Perplexity klasyfikacji: $\exp(\mathrm{CE}_\mathrm{test})$. Klasyfikator zgadujący równomiernie ma
        wartość 3; mniej oznacza pewne i trafne predykcje, więcej --- pewne i błędne.
  \item Perplexity LM: średnia NLL następnego tokenu na tekstach testowych (tylko prawdziwe tokeny, bez
        dopełnienia), liczona głowicą LM backbone'u (\texttt{lm\_head} miniLLaMA, a dla GPT-2 wiązane
        embeddingi \texttt{wte}) przed fine-tunem i po nim. Ponieważ tokenizery różnią się (BPE 32k vs GPT-2 50k),
        porównania między rodzinami modeli robimy w \emph{bitach na znak}.
\end{itemize}

\section{Wyniki liczbowe}
Wszystkie liczby dotyczą oficjalnego splitu testowego (12\,284 tweety), dla wag z epoki o najlepszym
val F1 macro. Wyniki pochodzą z pojedynczego uruchomienia (jedno ziarno), więc różnice rzędu
0,5\,pp traktujemy jako mieszczące się w szumie.

\subsection{Ranking wszystkich wariantów}
\maybeinput{generated/results_table.tex}

\begin{figure}[H]
\centering
\maybefig[]{generated/figures/accuracy_ranking.pdf}
\caption{Accuracy na teście dla całej siatki (pełny zbiór treningowy); kolor = backbone.}
\label{fig:acc-ranking}
\end{figure}

\begin{figure}[H]
\centering
\maybefig[]{generated/figures/f1_ranking.pdf}
\caption{F1 macro na teście dla całej siatki.}
\label{fig:f1-ranking}
\end{figure}

\subsection{Metoda dostrajania a backbone}
\begin{figure}[H]
\centering
\maybefig[]{generated/figures/method_heatmap.pdf}
\caption{Accuracy [\%] w układzie backbone $\times$ metoda. Kolumny ,,top $p$'' to odmrożenie
$k=\lceil pL\rceil$ ostatnich bloków; dla płytkich modeli kilka udziałów daje to samo $k$, wtedy
ten sam run wypełnia kilka kolumn.}
\label{fig:heatmap}
\end{figure}

\subsection{Rozmiar modelu}
\begin{figure}[H]
\centering
\maybefig[]{generated/figures/acc_vs_size.pdf}
\caption{Accuracy w funkcji rozmiaru backbone'u (parametry modelu bazowego, bez adapterów).
Dla ,,ostatnich bloków'' pokazany najlepszy wariant $k$.}
\label{fig:size}
\end{figure}

\subsection{Koszt treningu}
\begin{figure}[H]
\centering
\maybefig[]{generated/figures/tradeoff.pdf}
\caption{Accuracy vs liczba trenowanych parametrów (lewo) i vs czas treningu 5 epok na 2$\times$RTX 3060 (prawo).}
\label{fig:tradeoff}
\end{figure}

\subsection{Perplexity}
\maybeinput{generated/lm_table.tex}

\begin{figure}[H]
\centering
\maybefig[]{generated/figures/perplexity.pdf}
\caption{Lewo: perplexity klasyfikacji $\exp(\mathrm{CE})$ (linia przerywana = zgadywanie, 3).
Środek: perplexity LM na token przed i po fine-tunie (skala log; tokenizery różnią się między
rodzinami). Prawo: bity na znak --- porównywalne między tokenizerami.}
\label{fig:ppl}
\end{figure}

\subsection{Inferencja}
\begin{figure}[H]
\centering
\maybefig[]{generated/figures/inference.pdf}
\caption{Przepustowość (cały test, batch 32) i opóźnienie (batch 1), jedna karta RTX 3060.}
\label{fig:inference}
\end{figure}

\subsection{Ablacja rozmiaru danych}
\maybeinput{generated/ablation_table.tex}

\begin{figure}[H]
\centering
\maybefig[]{generated/figures/data_ablation.pdf}
\caption{Accuracy w funkcji liczby przykładów treningowych (2\,000, 10\,000, pełny split).}
\label{fig:ablation}
\end{figure}

\subsection{Przebieg treningu}
\begin{figure}[H]
\centering
\maybefig[]{generated/figures/learning_curves.pdf}
\caption{Krzywe uczenia dla wybranych runów: strata treningowa, strata walidacyjna oraz accuracy
(walidacja --- pełne linie z punktami, trening --- blade).}
\label{fig:curves}
\end{figure}

\begin{figure}[H]
\centering
\maybefig[]{generated/figures/val_f1_curves.pdf}
\caption{Val F1 macro po każdej epoce dla pełnego fine-tune'u i LoRA.}
\label{fig:valf1}
\end{figure}

\subsection{Macierze pomyłek}
\begin{figure}[H]
\centering
\maybefig[]{generated/figures/confusion_best.pdf}
\caption{Macierze pomyłek na teście dla najlepszego własnego modelu i najlepszego wariantu GPT-2
(liczności i procent wiersza).}
\label{fig:cm}
\end{figure}

\subsection{Status kolejki}
\maybeinput{generated/missing_runs.tex}

\section{Analiza porównawcza}

\subsection{Pretrening vs trening od zera}
Przy identycznej architekturze (miniLLaMA-5M, pełny fine-tune) start z checkpointu LM daje
\resptFiveFullAcc\,\% accuracy wobec \resscratchAcc\,\% od zera (+\resgainPretrain\,pp; F1 macro 61,1
vs 55,2). Przewaga utrzymuje się na każdym rozmiarze danych (tab.~\ref{tab:ablation}): pt-110M
wyprzedza scratch o 5,7\,pp przy 2\,000 przykładów, 8,4\,pp przy 10\,000 i 6,2\,pp na pełnym zbiorze.
Model od zera trenowany na całych 45\,615 tweetach osiąga mniej więcej poziom pt-110M uczonego na
10\,000 przykładów (57,7 vs 58,5\,\%), a GPT-2 z LoRA na 10\,000 przykładów (68,9\,\%) bije każdy
własny model trenowany na pełnych danych. Pretrening działa więc jak mnożnik danych: wiedza
o języku nie musi być odkrywana z etykiet sentymentu.

Jednocześnie pretrening na WikiText-103 nie daje cech, które same w sobie kodują sentyment:
zamrożony backbone z wytrenowaną tylko głową (48,7--53,9\,\%) przegrywa ze scratch-5M, a pt-5M
z jednym odmrożonym blokiem (56,5\,\%) jest nadal nieco poniżej. Korzyść z pretreningu ujawnia się
dopiero wtedy, gdy reprezentacje mogą się dopasować do zadania.

\textbf{Kiedy trening od zera ma sens?} W naszych wynikach scratch nigdy nie wygrywa z pretrenowanym
modelem o tym samym rozmiarze i tej samej liczbie trenowanych parametrów, a koszt inferencji jest
identyczny (ta sama architektura). Trening od zera jest uzasadniony, gdy nie istnieje checkpoint
pasujący do domeny lub tokenizacji (np. nietypowy alfabet, sekwencje niebędące tekstem), gdy
licencja lub prywatność wyklucza użycie cudzych wag, albo gdy danych etykietowanych jest bardzo
dużo, a domena pretreningu jest odległa --- wtedy przewaga pretreningu maleje, a pełna kontrola
nad słownikiem i architekturą może być ważniejsza.

\subsection{Wpływ rozmiaru modelu}
Dla własnych modeli accuracy pełnego fine-tune'u rośnie tylko nieznacznie: 62,7 (5M), 62,7 (10M),
63,6 (25M), 63,9 (50M), 64,0\,\% (110M), czyli plateau od ok.\ 25M (rys.~\ref{fig:size}). Przyczyną
jest pretrening, a nie pojemność: 50M i 110M zostały przerwane limitem 90\,min po ok.\ 2,25 i 1,3
epoki, a perplexity LM na tweetach przed fine-tunem jest dla 25M, 50M i 110M praktycznie taka sama
(ok.\ 3,2 bitu na znak). Większe modele nie nauczyły się więcej o języku, więc nie mają czego
przenieść. Silniej od rozmiaru zależą warianty z zamrożonym backbone'em: sama głowa rośnie z 48,7
do 53,9\,\%, a LoRA z ok.\ 60\,\% (5M, 10M) do 63,3--63,9\,\% (od 25M).

Najmocniejszym argumentem jest porównanie GPT-2 z pt-110M: oba mają 12 bloków i $d=768$, a różnica
na pełnym fine-tunie wynosi 6,4\,pp (70,4 vs 64,0\,\%). Przed fine-tunem GPT-2 modeluje tweety
z 1,84 bitu na znak, a pt-110M z 3,21. GPT-2 uczył się na WebText (ok.\ 40\,GB tekstu z internetu,
w tym języka nieformalnego), a nasze modele na ok.\ 100\,M tokenów encyklopedycznego WikiText-103.
O jakości transferu decyduje zatem skala i dopasowanie domenowe danych pretreningu, a nie liczba
parametrów. Nawet GPT-2 z samą głową (62,0\,\%) prawie dorównuje pełnemu fine-tune'owi pt-5M
i pt-10M (62,7\,\%).

\subsection{Porównanie metod dostrajania}
\begin{itemize}
  \item \textbf{Full i LoRA} są na szczycie dla każdego backbone'u. Od 25M w górę różnica między nimi
        nie przekracza 0,35\,pp (dla pt-50M jest zerowa, dla GPT-2 LoRA wygrywa o 0,3\,pp:
        \resbestAllAcc\,\% / F1 \resbestAllFone). Dla najmniejszych modeli (5M, 10M) LoRA traci
        2,2--2,7\,pp: adaptery rzędu 16 w modelach o $d=144$ i $240$ to mało swobody, a tak małe
        reprezentacje wymagają głębszej przebudowy.
  \item \textbf{Ostatnie $k$ bloków}: accuracy rośnie prawie monotonicznie z $k$. Przy odmrożeniu
        ok.\ połowy bloków strata do full wynosi 0,6\,pp dla GPT-2 (69,8 vs 70,4\,\%), 1,2--2,2\,pp dla
        miniLLaMA od 25M, ale 2,6\,pp dla 10M i 6,1\,pp dla 5M, gdzie ,,połowa'' to jeden z dwóch
        bloków. W GPT-2 już 2 z 12 bloków dają 67,8\,\%, więcej niż jakikolwiek własny model.
  \item \textbf{Sama głowa} jest wyraźnie najgorsza: 10--14\,pp za full dla miniLLaMA (strata
        maleje z rozmiarem) i 8,4\,pp dla GPT-2. Cechy ostatniego tokena modelu LM nie są
        separowalne względem sentymentu przez małe MLP.
\end{itemize}

\subsection{Koszt: parametry, czas, inferencja}
LoRA w GPT-2 trenuje 1,9\,\% wag i daje najlepszy wynik przy 38\,\% krótszym treningu niż full
(10,3 vs 16,6\,min). Odmrożenie połowy bloków to kolejne 20\,\% oszczędności (8,4\,min) kosztem 0,6\,pp.
W małych miniLLaMA LoRA \emph{nie} skraca treningu (pt-5M: 3,2 vs 2,6\,min, pt-10M: 4,2 vs
3,6\,min): oszczędność na gradientach wag i stanie optymalizatora jest mniejsza niż narzut
dodatkowych mnożeń i dropoutu adapterów w forwardzie. Zysk z LoRA pojawia się dopiero przy
większych macierzach (pt-110M: 14,0 vs 16,1\,min).

Przepustowość inferencji zależy praktycznie tylko od rozmiaru backbone'u: od ok.\ 5\,800 przykł./s
(5M) do ok.\ 540 (110M) i 620 (GPT-2). Wyjątkiem są niezłączone adaptery LoRA, które zwiększają
opóźnienie pojedynczego przykładu: GPT-2 25,6 vs 13,5\,ms, pt-110M 37 vs 24\,ms. Przy wdrożeniu
należy wcześniej złączyć adaptery z wagami (\texttt{merge\_and\_unload}); wtedy koszt inferencji jest
taki jak dla full. GPT-2 jest szybszy od pt-110M, choć ich kształt i koszt FFN na token są
porównywalne (SwiGLU $3\times768\times2048$ vs MLP $2\times768\times3072$). Różnica wynika najpewniej
z implementacji: HF używa zoptymalizowanych kerneli, a miniLLaMA wykonuje dodatkowo RoPE
i trzy mnożenia w FFN.

\subsection{Perplexity}
\textbf{Klasyfikacja.} Wszystkie runy mają $\exp(\mathrm{CE}) < 3$ (lepiej niż zgadywanie):
od ok.\ 1,9 (GPT-2 full) do 2,8 (sama głowa). Niska perplexity klasyfikacji idzie w parze
z wysokim accuracy, ale nie jest z nim tożsama: GPT-2 full ma najniższą stratę, a LoRA najwyższe
accuracy. Kalibrację pogarsza przeuczenie, o czym niżej.

\textbf{Modelowanie języka.} Przed fine-tunem GPT-2 ma na tweetach perplexity 182 (1,84\,bpc),
a własne modele od ok.\ 1\,800 do 4\,800 na token (3,19--3,60\,bpc). To skutek przesunięcia
domeny (Wikipedia $\rightarrow$ tweety) i mniejszego korpusu. Po fine-tunie perplexity LM rośnie:
dla GPT-2 full 8,5-krotnie (do ok.\ 1\,540), dla LoRA 3,4-krotnie (do ok.\ 620), a dla miniLLaMA
full ok.\ 1,8--2,7-krotnie. Ukryte stany przesuwają się w stronę tego, co potrzebne głowie
klasyfikacyjnej, a niewytrenowana głowica LM przestaje do nich pasować, czyli zachodzi
\emph{zapominanie} zdolności LM. LoRA zapomina wyraźnie mniej niż full, bo zmiana wag jest niskiego
rzędu. W GPT-2 jeszcze mniej zapomina odmrożenie ostatnich bloków: wzrost od 1,15-krotnego (2 bloki,
209) do 2,4-krotnego (6 bloków, 440), bo dolna część sieci zostaje nietknięta.
Sama głowa zostawia perplexity bez zmian, co jest jednocześnie testem poprawności pomiaru
(backbone zamrożony).

\subsection{Zbieżność, stabilność i przeuczenie}
Najlepsza epoka walidacyjna to najczęściej 2 lub 3 (23 z 37 runów siatki). Przy pełnym fine-tunie
strata walidacyjna jest minimalna po 1--2 epokach, a potem rośnie z ok.\ 0,7 do 1,2--1,7, podczas
gdy accuracy treningowe dochodzi do 95--99,9\,\% (rys.~\ref{fig:curves}). Accuracy walidacyjne
spada wolniej niż rośnie strata: model coraz pewniej się myli. Bez early stoppingu kluczowy jest
wybór \texttt{best\_eval}: pełny fine-tune osiągał maksimum w epoce 1--2, LoRA w epoce 2--3
(z wyjątkiem pt-10M --- epoka 4).
LoRA zachowuje się stabilniej: dla GPT-2 strata walidacyjna zostaje w zakresie 0,62--0,67 przez
wszystkie 5 epok, a strata treningowa nie spada poniżej 0,4, więc niski rząd działa jak
regularyzacja. Warianty o małej pojemności uczą się wolniej: sama głowa osiąga maksimum w epoce
4--5, a GPT-2 z 2--4 odmrożonymi blokami w epoce 5, więc tu dłuższy trening mógłby jeszcze pomóc.

Accuracy testowe jest systematycznie niższe od walidacyjnego (np. pt-110M: val F1 ok.\ 68 vs test
63). Rozkład klas w teście różni się od treningowego: negatywnych jest dwukrotnie więcej (32 vs
16\,\%), a pozytywnych dwukrotnie mniej (19 vs 39\,\%). Widać to w macierzach pomyłek
(rys.~\ref{fig:cm}): najwięcej błędów to tweety negatywne i pozytywne klasyfikowane jako neutralne.
GPT-2 rozpoznaje klasę negatywną znacznie lepiej (recall 75 vs 59\,\%), co tłumaczy większość jego
przewagi.

\subsection{Spostrzeżenia praktyczne}
\begin{itemize}
  \item \textbf{Pamięć hosta, a nie GPU, była wąskim gardłem.} Na końcu \texttt{fit} w trybie
        \texttt{ddp\_spawn} w ciągu kilku sekund zapisywano i wczytywano kilka kopii wag (checkpoint
        epoki, \texttt{final.pt}, tymczasowy checkpoint Lightninga, \texttt{best\_eval}). Przy 500\,MB
        na kopię GPT-2 to wystarczyło, żeby zawiesić WSL2. Pomogło usunięcie zbędnego zapisu,
        zwolnienie loaderów i trenera przed testem oraz logowanie pamięci
        (\texttt{mem\_log.csv} w każdym runie).
  \item \textbf{Perplexity LM łatwo policzyć źle.} Pierwsza wersja wliczała do sumy NLL pozycje
        dopełnienia, a dzieliła tylko przez prawdziwe tokeny; wartości przekraczały próg $e^{20}$
        i przed/po wychodziło identyczne. Po zamaskowaniu celów na dopełnieniu wyniki stały się
        sensowne (np.\ 182 dla GPT-2, zgodnie z przesunięciem domeny).
  \item \textbf{PEFT \texttt{modules\_to\_save} dopasowuje nazwy po sufiksie}: wpis \texttt{"head"}
        objął także \texttt{backbone.lm\_head}, który dostał nieużywaną trenowalną kopię
        (vocab$\times d$). Nie wpływa to na wyniki klasyfikacji (warstwa nie jest na ścieżce forwardu
        i nie dostaje gradientu), ale zawyża licznik trenowanych parametrów, więc w tabelach jest
        on skorygowany. Faktycznie LoRA w miniLLaMA trenuje 2,2--2,7\,\% wag.
  \item \textbf{Checkpoint po każdej epoce} kosztuje dużo dysku (do 5$\times$0,5\,GB na run), ale
        pozwolił wybrać najlepszą epokę post hoc i przeliczyć metryki LM bez ponownego treningu.
\end{itemize}

\section{Wnioski}
\begin{enumerate}
  \item Dostrajanie pretrenowanego modelu jest lepsze od treningu od zera przy każdym badanym
        rozmiarze danych. Ten sam 5M model zyskuje \resgainPretrain\,pp accuracy, a pretrenowane
        modele osiągają z 10\,000 przykładów wynik, którego scratch nie osiąga z 45\,000.
  \item Jakość i domena danych pretreningu ważą więcej niż liczba parametrów: GPT-2 (\resbestGptAcc\,\%)
        wyprzedza najlepszy własny model (\resbestPtLabel, \resbestPtAcc\,\%) o \resgapPtGpt\,pp przy
        podobnym rozmiarze, a własne modele saturują od ok.\ 25M, bo większe były niedotrenowane.
  \item LoRA to najlepszy kompromis dla modeli od ok.\ 25M parametrów: jakość full przy 2--3\,\%
        trenowanych wag, krótszym treningu (dla GPT-2 o 38\,\%) i mniejszym zapominaniu LM.
        W bardzo małych modelach lepszy jest pełny fine-tune.
  \item Odmrożenie połowy ostatnich bloków to tańsza alternatywa dla głębszych modeli (strata
        0,6--2,2\,pp do full), ale nie dla płytkich. Trening samej głowy na zamrożonym backbone'ie jest
        niewystarczający ($-$8 do $-$14\,pp).
  \item Pełny fine-tune i LoRA przeuczają się po 2--3 epokach. Wybór najlepszej epoki po val F1 jest
        niezbędny, a dla tych metod 3 epoki zamiast 5 skróciłyby trening o 40\,\% bez straty jakości.
        Warianty z zamrożoną większością wag potrzebują więcej epok.
  \item Do wdrożenia: złączyć adaptery LoRA (inaczej opóźnienie rośnie ok.\ dwukrotnie) i pamiętać,
        że zmiana rozkładu klas między treningiem a testem obniża wynik bardziej niż różnice
        między metodami dostrajania.
\end{enumerate}

\appendix
\section{Reprodukcja}
\begin{verbatim}
cd lab03
./background_process_start.sh  # kolejka fine-tune (log: nohup.out)
./shell.sh uv run python src/recompute_lm_eval.py  # metryki LM
./report/build.sh              # generator raportu + pdflatex
./shell.sh uv run python src/finetune.py --list-sources
\end{verbatim}

\end{document}